\section{The original tensor windows and the complete boundary--kernel sum}
\label{sec:twa-original-tensor-windows}

This calculation uses the complete proved TW1--TW21 and FPK.1--FPK.29
bodies retained with this audit. Fix the original full-order hypothetical
quartet of BRU1, with $0<\delta<1/2$, $\gamma>2$, and positive integer
order $m$. The original $h$, $v_h=2\xi/h$, its full Taylor unit and
the coefficient observation $\Lambda$ of BRU3 remain unchanged.
Only the positive integer tensor order $k$ varies. For every such $k$ set
\[
\begin{gathered}
\ell_k=1+k(m-1),\quad q=\ell_k(k+1)^2,\quad c=k/2,\\
\lambda_{a,b}=c+(2a-k)\delta+i(2b-k)\gamma,
\quad 0\le a,b\le k,\\
\chi(S)=\prod_{a,b=0}^k(S-\lambda_{a,b})^{\ell_k},
\quad E=\mathbb C[S]/(\chi),\qquad S=c+iu,\\
w_h(u)=|v_h(1/2+iu)|^2/(2\pi),\quad
m_{h,k}=w_h^{*k},\quad \alpha=\pi/2,\\
\vartheta_h=\int_{-1}^1w_h(u)\,du,\qquad
M_\alpha=\int_{\mathbb R}e^{\alpha u}w_h(u)\,du.
\end{gathered}\tag{TWA1}
\]
All integrations retain these measures and their masses. In particular
no probability measure is substituted for $w_h$, $m_{h,k}$ or a Gamma
comparison. The moment and source estimates below concern $k\ge3$;
the algebraic BRU decomposition itself retains its domain $k\ge1$.

\subsection{The complete proved dependence of the comparison constants}

Retain precisely the constants of TW2--TW12:
\[
\begin{gathered}
p=8m-\tfrac92,\quad B=42+8m,\quad M=B/2=21+4m,\\
c_\Gamma=\sqrt2e^{-7/6},\qquad
C_\Gamma=\sqrt{2\sqrt5}\,e^{2/3},\qquad
\sigma(u)=|\Gamma(1/4+iu/2)|^2/(2\pi),\\
a_k=\frac{c_h\vartheta_h^{k-3}e^{-\alpha(k-3)}(k-2)^{-B}}
{C_\Gamma2^{B/2}},\qquad
A_k=\frac{C_h^{\rm tilt}}{c_\Gamma}k^{p+1}M_\alpha^{k-1},\\
D_N=[1+4(N+M)^2]^M.
\end{gathered}\tag{TWA2}
\]
Here $c_h>0$ is the original proved convolution lower-envelope
constant of TW7, and $C_h^{\rm tilt}$ is exactly the maximum of the
compact and tail expressions BT12. These depend on the fixed actual
quartet, not on $k$ or $N$. The complete BT proof retains the original
zeta numerator and proves both $0<M_\alpha<\infty$ and
$e^{\alpha|u|}w_h(u)\le C_h^{\rm tilt}(1+|u|)^{-p}$.
The source exponent at the exponential boundary only gives the
finite derivative orders $j<p-1$ there; it is not a claim of all
boundary derivatives. The un-tilted polynomial moments used here are
all finite because the original exponential factor is retained.

The full-polynomial inequality TW12, transported by the exact algebra
isomorphism $P(S)\mapsto P(c+iu)$ and its inverse
$f(u)\mapsto f((S-c)/i)$, is
\[
\frac{a_k}{D_N}\int|P(c+iu)|^2\sigma(u)\,du
\le\int|P(c+iu)|^2m_{h,k}(u)\,du
\le A_k\int|P(c+iu)|^2\sigma(u)\,du,
\quad\deg P\le N.\tag{TWA3}
\]
The map retains the original complex coordinate and the original
remainder polynomial. TW proves this for every complex coefficient
vector, so it applies to every vector in each original affine
remainder fibre as well as to every monic polynomial.

Define two actual constants, with no change to either measure,
\[
\begin{aligned}
\mathfrak s_h&=\alpha+\log(M_\alpha/\vartheta_h),\\
\mathfrak c_h&=-\log M_\alpha+3\log\vartheta_h-3\alpha
 +\log\frac{C_h^{\rm tilt}C_\Gamma2^{B/2}}{c_\Gamma c_h}.
\end{aligned}
\]
Subtraction of the logarithms of the two constants in TWA2 gives
the exact identity
\[
\log(A_k/a_k)=k\mathfrak s_h+(p+1)\log k
                      +B\log(k-2)+\mathfrak c_h.\tag{TWA4}
\]
Moreover $\mathfrak s_h>\alpha$. Indeed the original evenness gives
$M_\alpha=\int\cosh(\alpha u)w_h(u)\,du$.
Its restriction to $[-1,1]$ is at least $\vartheta_h$ because
$\cosh(\alpha u)\ge1$. Its complement has strictly positive integral:
$w_h$ is positive almost everywhere, and that complement has positive
measure. Thus $M_\alpha>\vartheta_h>0$, proving the strict assertion.
This evaluates the sign of the actual source constant; it does not
replace either of its defining integrals.

\subsection{A finite two-sided error for the literal monic windows}

Let $\omega_{h,k,n}$ be the original minimum squared norm of a monic
degree-$n$ polynomial in $S$. Let
\[
\gamma_n=\sqrt{2\pi}\frac{(2n)!}{4^n}
\]
be the squared norm for the fixed one-factor Gamma comparison.
The degree-specific map $P\mapsto i^{-n}P(c+iu)$ is a bijection of
the two monic affine spaces and preserves the norm because $|i|=1$.
It is used only on that stated affine space, not as a new algebra map.
Minimization of TWA3 therefore gives exactly TW14:
\[
a_k\gamma_n/D_n\le\omega_{h,k,n}\le A_k\gamma_n.
\]
Divide these two positive bounds at each of the two original
endpoint pairs. Define
\[
\begin{aligned}
W_k&=\log\frac{\omega_{h,k,2q-1}\omega_{h,k,2q}}
{\omega_{h,k,q-1}\omega_{h,k,q}},\\
W_k^\sigma&=\log\frac{\gamma_{2q-1}\gamma_{2q}}
{\gamma_{q-1}\gamma_q}
 =\log\frac{(4q-2)!(4q)!}{4^{2q}(2q-2)!(2q)!},\\
E_k^+&=2\log(A_k/a_k)+\log D_{q-1}+\log D_q,\\
E_k^-&=2\log(A_k/a_k)+\log D_{2q-1}+\log D_{2q},\qquad
\tau_k^{\rm win}=W_k-W_k^\sigma.
\end{aligned}\tag{TWA5}
\]
The symbol $\tau_k^{\rm win}$ denotes this scalar window error;
it does not rename the original phase polynomial $\sum_i\Phi_h(s_i)$.
Every comparison denominator is positive. Applying TWA3 to any
nonzero degree-$n$ polynomial gives $A_kD_n/a_k\ge1$, so
$E_k^+,E_k^-\ge0$. Dividing the lower bound for each numerator by
the upper bound for its denominator, and doing the reverse for the
upper bound, proves
\[
\begin{gathered}
 -e_k^-\le\tau_k^{\rm win}\le e_k^+,\qquad
 e_k^-=\min\{E_k^-,\Xi_{k,2q}-\Xi_{k,2q-2}\},\\
 e_k^+=\min\{E_k^+,\Xi_{k,q}-\Xi_{k,q-2}\},
 \qquad 0\leq e_k^\pm\leq E_k^\pm .
\end{gathered}\tag{TWA6}
\]
Here the exact source determinant deficit is
$\Xi_{k,N}=(N+1)\log A_k-
\log(\det H_{k,N}/\det H_N^{\sigma,k})$ from HMR36,
with its full finite bound HMR38 and complete DMR proof.
This use also requires the intermediate source degrees
$q-2$ and $2q-2$, where that finite determinant is defined.
To prove the new intersection, let $\Xi_{k,-1}=0$ for
the empty source. The monic orthogonal change of frame has
diagonal one, so
$\det H_{k,n}/\det H_{k,n-1}=\omega_{h,k,n}$
and
$\det H_n^{\sigma,k}/\det H_{n-1}^{\sigma,k}=\gamma_n$.
It follows exactly that
$z_{k,n}:=\Xi_{k,n}-\Xi_{k,n-1}
=\log(A_k\gamma_n/\omega_{h,k,n})\geq0$,
where positivity is the upper monic inequality just proved.
Taking the same four norm logarithms as in TWA5 gives
\[
 \tau_k^{\rm win}
 =(z_{k,q-1}+z_{k,q})-(z_{k,2q-1}+z_{k,2q}).
\]
The two sums telescope to the respective deficit differences
in TWA6. Dropping either nonnegative sum gives the additional
two-sided bound. Intersecting it with the previously proved
$[-E_k^-,E_k^+]$ proves TWA6. Thus every individual monic
norm and both adjacent differences are retained.

All constants in TWA4 are fixed as $k$ varies. Since
$\log D_N=M\log(1+4(N+M)^2)=O_h(1+\log(N+1))$, TWA4 proves
$E_k^\pm=O_h(k+\log q)$ on the stated growing windows. The exact
polynomial
\[
q=(m-1)k^3+(2m-1)k^2+(m+1)k+1
\]
gives $q\ge(k+1)^2$, $k/q\to0$ and $\log q/q\to0$. Consequently
$E_k^\pm/q\to0$. The same exact identities give
$E_k^\pm/(2k)\to\mathfrak s_h>\alpha$; in particular these errors
are small after division by $q$, but not after division by $k$.

For completeness, the exact fixed Gamma window has the finite bound
\[
4q\log(4q/e)-\log6\le W_k^\sigma\le4q\log(4q/e).
\tag{TWA7}
\]
To verify it, write
$\log(\gamma_{2q}/\gamma_q)=2q\log(4q/e)+\eta_q$.
The sum $\sum_{j=2q+1}^{4q}\log j$ exceeds
$\int_{2q}^{4q}\log x\,dx$ by a number in $[0,\log2]$, because
each summand is between the adjacent lower and upper integral sums.
Thus $0\le\eta_q\le\log2$. Direct factorial division also gives
\[
\frac{\gamma_{2q-1}}{\gamma_{q-1}}
=\frac{\gamma_{2q}}{\gamma_q}\frac{q-1/2}{4q-1},
\qquad 1/6\le(q-1/2)/(4q-1)<1/4.
\]
Adding logarithms proves TWA7. The exact factorial expression TWA5
is retained whenever that sharper finite value is used.

\subsection{The complete BRU20 sum, with every original cost}

Use the original arithmetic minimum sections at $x=1$, the actual
kernel and boundary of $\Lambda$, and precisely the original row
lengths $a_j(1)$ and contractions $\beta_j(1),\nu_j(1)$ of BRU8.
These row lengths are distinct from the comparison constant $a_k$
in TWA2. The two scalar sums remain
\[
\begin{aligned}
F_K(1)&=\sum_{j=q-1}^{2q-1}\omega_j
          \log(1+\beta_j(1)/a_j(1)),\\
F_B(1)&=\sum_{j=q-1}^{2q-1}\omega_j
          \log\left(1+\frac{\nu_j(1)}{a_j(1)+\beta_j(1)}\right),\\
\omega_{q-1}&=\omega_{2q-1}=1,
\qquad\omega_j=2\quad(q\le j\le2q-2).
\end{aligned}\tag{TWA8}
\]
The symbol $\omega_j$ in this formula denotes only its displayed
endpoint weight; the original monic squared norms retain all three
indices $\omega_{h,k,n}$ in TWA5. BRU11 proves the exact determinant
product for each summand and BRU21 proves
$F_K(1)+F_B(1)=\mathcal B_k$. No value of the unitful observation or
its kernel is changed in that identity.

Retain the original HC phase $\phi_N$, relative control
$\varepsilon_N$, contraction penalties $P_-,P_+$ and full trace
\[
\mathcal L_k=2\delta\ell_k(k+1)
                  \left\lfloor\frac{(k+1)^2}{4}\right\rfloor.
\]
Define the following sum of the original nonnegative terms:
\[
\begin{aligned}
\mathcal S_k={}&P_-+P_+
 +\sum_{N=q-1}^{2q-1}\omega_N
                   \log(1+\phi_N^2/\mathcal L_k^2)\\
&+\sum_{N=q-1}^{2q-1}\omega_N
       \log\frac{\varepsilon_N^2+\phi_N^2}
                    {\mathcal L_k^2+\phi_N^2}.
\end{aligned}\tag{TWA9}
\]
The original proofs give $0<d_j<1$, so each term of $P_-,P_+$
is nonnegative, and give $\varepsilon_N\ge\mathcal L_k>0$ on all
displayed degrees. Thus $\mathcal S_k\ge0$. Its entire values,
phases and multiplicities remain in TWA9.

Since $\sum_N\omega_N=2q$, splitting each positive
$\varepsilon_N^2+\phi_N^2$ into its two displayed factors gives
exactly
\[
\boxed{F_K(1)+F_B(1)
 =4q\log\mathcal L_k-W_k^\sigma
                      +\mathcal S_k-\tau_k^{\rm win}.}
\tag{TWA10}
\]
Indeed HC2 is $\mathcal B_k=\sum_N\omega_N
\log(\varepsilon_N^2+\phi_N^2)-W_k+P_-+P_+$.
For each summand,
\[
\log(\varepsilon_N^2+\phi_N^2)
=2\log\mathcal L_k+\log(1+\phi_N^2/\mathcal L_k^2)
 +\log\frac{\varepsilon_N^2+\phi_N^2}
                   {\mathcal L_k^2+\phi_N^2}.
\]
Insert this identity, TWA5, and BRU21 to prove TWA10 with no
discarded term. TWA6 then gives the full finite interval
\[
\begin{aligned}
4q\log\mathcal L_k-W_k^\sigma+\mathcal S_k-e_k^+
&\le F_K(1)+F_B(1)\\
&\le4q\log\mathcal L_k-W_k^\sigma+\mathcal S_k+e_k^-.
\end{aligned}\tag{TWA11}
\]
This upper inequality retains $\mathcal S_k$; no bound on its
growing value follows by declaring its terms nonnegative.

Using the proved upper side of TWA7 in the lower side of TWA11 gives
\[
\begin{aligned}
F_K(1)+F_B(1)\ge{}&4q\log(\delta e k/8)-e_k^++\mathcal S_k\\
&+4q\log\frac{2\mathcal L_k}{\delta kq}.
\end{aligned}\tag{TWA12}
\]
The identity behind its constant is
$\mathcal L_k e/(4q)=(\delta e k/8)
[2\mathcal L_k/(\delta kq)]$. The last factor is retained exactly:
\[
\frac{2\mathcal L_k}{\delta kq}
=\begin{cases}(k+2)/(k+1),&k\text{ even},\\
(k+1)/k,&k\text{ odd}.
\end{cases}
\]
It is larger than one and tends to one. Therefore TWA6, TWA9 and
TWA12 prove the valid growing-order conclusion
\[
\boxed{\liminf_{k\to\infty}
\frac{F_K(1)+F_B(1)-4q\log k}{q}
\ge4\log(\delta e/8).}\tag{TWA13}
\]
Every limit here keeps the actual original quartet fixed. This
lower bound is on the combined original sum. At boundary rank zero
BRU gives $F_B=0$, and at full boundary rank it gives $F_K=0$;
TWA10--TWA13 retain both cases without assigning either separate
term a positive lower bound at intermediate rank.

\subsection{Exact finite comparison with the full reference quotient}

Let $G_N^\sigma$ denote the quotient Gram of the same remainder
map $J_N:\mathcal P_N\to E$ for the fixed measure $\sigma$ on
$S=c+iu$. The source identity on $\mathcal P_N$ commutes with
$J_N$ and the complete original observation. Minimizing both sides
of TWA3 over its identical affine fibre $J_NP=z$ proves
\[
\max\{a_k/D_N,A_ke^{-\Xi_{k,N}}\}G_N^\sigma
       \preceq G_N(1)\preceq A_kG_N^\sigma.
\tag{TWA14}
\]
The additional lower factor follows before quotienting:
the eigenvalues of $H_N^{\sigma,k}{}^{-1/2}H_{k,N}
H_N^{\sigma,k}{}^{-1/2}/A_k$ lie in $(0,1]$ and have product
$e^{-\Xi_{k,N}}$, so every eigenvalue is at least that product.
Minimizing the resulting full-vector bound over the same
remainder fibre gives this factor. Taking the larger of two
proved lower constants proves the displayed maximum.
For a positive matrix inequality $lG\preceq H\preceq uG$,
congruence with $G^{-1/2}$ bounds each of the $q$ eigenvalues
between $l$ and $u$, and hence bounds its determinant between
$l^q\det G$ and $u^q\det G$. Apply this argument to TWA14
at the four original endpoints, keeping their signs. Define
\[
F_{\sigma,k}=\log\frac{\det G_{q-1}^\sigma\det G_q^\sigma}
{\det G_{2q-1}^\sigma\det G_{2q}^\sigma},\qquad
\delta_k^\sigma=F_K(1)+F_B(1)-F_{\sigma,k}.
\]
The resulting full finite enclosure is
\[
\begin{gathered}
 \underline\delta_k^\sigma\leq\delta_k^\sigma
                         \leq\overline\delta_k^\sigma,\\
 \underline\delta_k^\sigma
 =\max\{-qE_k^+,-\Xi_{k,q-1}-\Xi_{k,q}\},\\
 \overline\delta_k^\sigma
 =\min\{qE_k^-,\Xi_{k,2q-1}+\Xi_{k,2q}\},\\
 |\delta_k^\sigma|\leq2\Xi_{k,2q},\qquad
 \frac{|\delta_k^\sigma|}{kq}
 =O_h\!\left[
 \left(\frac{\log(k+2)}{k}\right)^{(p-1)/p}\right].
\end{gathered}\tag{TWA15}
\]
Here is the determinant transport that proves the additional
interval, rather than treating a quotient as a restriction.
Set $b_N=q\log A_k-\log(\det G_N(1)/\det G_N^\sigma)$.
In the $H_N^{\sigma,k}$ isometric source coordinates,
split the positive contraction $D/A_k$ along the full kernel
of the original $J_N$ and its orthogonal complement.
Its Schur complement $T_N$ is the canonical quotient
relative Gram. Block elimination gives
$\det(D/A_k)=\det F_N\det T_N$, with both $F_N,T_N$
positive contractions. The original quotient coordinate
factor is exactly $G_N^\sigma{}^{1/2}$, so
$b_N=-\log\det T_N$ and
$0\leq b_N\leq\Xi_{k,N}$. The complete maps and ordered
block proof are HMR39--42; the empty relation determinant
is one at $N=q-1$.
Substitution at the original four endpoints yields
$\delta_k^\sigma=-b_{q-1}-b_q+b_{2q-1}+b_{2q}$.
This proves the additional interval.
Restriction of the original source by its literal coefficient
inclusion gives $\Xi_{k,N}\leq\Xi_{k,2q}$ for $N\leq2q$:
complete the associated Gamma-isometric injection to a
unitary frame and use the same positive Schur determinant.
Thus both endpoint sums are at most $2\Xi_{k,2q}$.
HMR38 and HMR44 prove the displayed rate with their actual
growing cutoff $R=(k/\log(k+2))^{1/p}$ and $q\geq(k+1)^2$.
The factor $q$ remains in the original coarse comparison;
it is the new Schur determinant calculation, with all
relation directions included, that improves its interval.
In particular the monic $e_k^\pm$ and the quotient deficit
in TWA15 are their separately proved exact quantities.

The common degree-$2q$ condition-number enclosure has
$\kappa_k=A_kD_{2q}/a_k\ge1$. Its logarithm obeys
\[
\log\kappa_k=k\mathfrak s_h+O_h(\log k+\log q).
\tag{TWA16}
\]
The full-spectrum AW comparison on this same coefficient source
has the inherited allowance $(2q-1)\log\kappa_k$. Its scale
relative to the unchanged trace is exact:
\[
\lim_{k\to\infty}\frac{(2q-1)\log\kappa_k}{\mathcal L_k}
=\frac{4\mathfrak s_h}{\delta}
>\frac{2\pi}{\delta}>4\pi.
\tag{TWA17}
\]
To prove the limit, divide TWA16 by $k$, and use
$\mathcal L_k/(qk)\to\delta/2$, which follows directly from
the two parity formulas following TWA12. The strict bounds use
$\mathfrak s_h>\alpha=\pi/2$ and $\delta<1/2$.
Thus this particular condition-number enclosure retains a
nonzero cost on the trace scale. The actual signed finite
comparison now has the stronger TWA15 estimate:
$\delta_k^\sigma/(kq)\to0$, and hence
$\delta_k^\sigma/\mathcal L_k\to0$ because
$\mathcal L_k/(kq)\to\delta/2>0$.
This conclusion uses its proved determinant deficit,
while TWA16--17 continue to state the exact value of the
earlier condition-number allowance.

\subsection{The full-jet asymptotic connected to every finite endpoint}

We now give an exact map from the FPK fixed-packet construction to
the present moving original packet. This supplies every finite
residual, rather than using its fixed-packet limit at an unproved
degree. Order the roots $\lambda_{a,b}$ and their normalized
Taylor derivatives $0\le j<\ell_k$ in the original order. Let
$\mathsf E_\chi$ be the full Taylor map of FPK.3, and retain
\[
(\mathsf E_\chi)_{(a,b,j),r}
=\begin{cases}\binom rj\lambda_{a,b}^{r-j},&r\ge j,\\0,&r<j.
\end{cases}
\]
The exact determinant is the retained confluent Vandermonde
$\prod_{i<a}(\lambda_a-\lambda_i)^{\ell_k^2}$.

For the fixed comparison let $p_n$ and $\gamma_n$ be FPK.2, and
form the full kernel
\[
(\mathsf K_{k,N})_{(i,j),(a,l)}
=\sum_{n=0}^N\frac{p_n^{[j]}(\lambda_i)
                         \overline{p_n^{[l]}(\lambda_a)}}{\gamma_n},
\quad N\ge q-1.
\tag{TWA18}
\]
The first $q$ monic polynomials span every remainder, so this
onto full-jet observation has positive Gram. Its minimum-source
quotient is exactly
\[
G_N^\sigma=\mathsf E_\chi^*\mathsf K_{k,N}^{-1}\mathsf E_\chi,
\quad
\det G_N^\sigma=|\det\mathsf E_\chi|^2/\det\mathsf K_{k,N}.
\tag{TWA19}
\]
These are finite equalities, valid with every nilpotent direction.

Put $w_i=\lambda_i-c$, and retain the complete functions
\[
C_\pm(w)=2^{-1/2\pm w/2}/\Gamma(1/4\pm w/2).
\]
On each off-critical block let $\epsilon_i=\operatorname{sgn}\Re w_i$,
and let $\mathsf T_{k,N,i}$ be multiplication on its full Taylor
jet by $[C_{\epsilon_i}(w)N^{\epsilon_iw/2}]^{-1}$.
The dominant Gamma argument has positive real part, so the
analytic germ and the multiplication map are invertible. On
each critical block let
$\mathsf T_{k,N,i}=\operatorname{diag}_{0\le j<\ell_k}
(\log N)^{-j-1/2}$. Since $N\ge q-1\ge3$, these maps exist at
every original endpoint. Let $\mathsf T_{k,N}$ be their ordered
direct sum. The off-critical limiting blocks are
\[
(\mathsf H_\epsilon)_{(i,j),(a,l)}
=\frac{(-\epsilon/2)^{j+l}(j+l)!}
{j!l!\,[\epsilon(w_i+\overline w_a)/2]^{j+l+1}},
\]
and the critical block at $w_i$ is
\[
(\mathsf H_i^0)_{j,l}
=\frac{|C_+(w_i)|^2}{2^{j+l}j!l!}\int_{-1}^1t^{j+l}\,dt.
\]
Their ordered block sum $\mathsf H_k$ is positive: the first
blocks are the Gram matrices of independent exponential
polynomials on $(0,1)$ proved in FPK.16, and the second are
positive monomial Grams retaining all jet degrees. No finite
kernel cross entry is removed in defining the following residual:
\[
\Delta_{k,N}^\sigma
=\mathsf H_k^{-1/2}
 (\mathsf T_{k,N}\mathsf K_{k,N}\mathsf T_{k,N}^*
                                      -\mathsf H_k)
 \mathsf H_k^{-1/2}.
\tag{TWA20}
\]
The exact identity $I+\Delta_{k,N}^\sigma
=\mathsf H_k^{-1/2}\mathsf T_{k,N}\mathsf K_{k,N}
\mathsf T_{k,N}^*\mathsf H_k^{-1/2}$ proves positivity at every
such degree. Its real logarithmic determinant is consequently
well-defined, whether or not the residual is small.

The exact two exponents for this quartet are
\[
\begin{aligned}
\mathcal A_\chi
 &=\delta\ell_k(k+1)\sum_{a=0}^k|2a-k|
 =\mathcal L_k,\\
\mathcal B_\chi
 &=\begin{cases}(k+1)\ell_k^2,&k\text{ even},\\0,&k\text{ odd}.
\end{cases}
\end{aligned}\tag{TWA21}
\]
For $k=2r$, the absolute-value sum is $2r(r+1)$; for $k=2r+1$
it is $2(r+1)^2$. Both are $2\lfloor(k+1)^2/4\rfloor$,
proving the first equality with the original trace constant.
A root has real part $c$ precisely when $a=k/2$, which exists
only for even $k$. It then contributes one block of length
$\ell_k$ for each of the $k+1$ imaginary indices, proving the second.

The full triangular determinants of $\mathsf T_{k,N}$ give
\[
|\det\mathsf T_{k,N}|^2
=N^{-\mathcal L_k}(\log N)^{-\mathcal B_\chi}
 \prod_{\Re w_i\ne0}|C_{\epsilon_i}(w_i)|^{-2\ell_k}.
\]
Taking determinants of TWA20 and then applying TWA19 proves
the exact identity
\[
\log\det G_N^\sigma
=-\mathcal L_k\log N-\mathcal B_\chi\log\log N
 +d_k^\sigma-\log\det(I+\Delta_{k,N}^\sigma),
\tag{TWA22}
\]
where
\[
d_k^\sigma=\log\frac{|\det\mathsf E_\chi|^2}
{(\det\mathsf H_k)
 \prod_{\Re w_i\ne0}|C_{\epsilon_i}(w_i)|^{2\ell_k}}.
\]
All constants, full jets and finite cross entries are retained.

Let $(N_0,N_1,N_2,N_3)=(q-1,q,2q-1,2q)$ and
$(\epsilon_0,\epsilon_1,\epsilon_2,\epsilon_3)=(1,1,-1,-1)$.
Define
\[
\begin{aligned}
\mathcal M_k={}&\mathcal L_k
 \log\frac{(2q-1)(2q)}{(q-1)q}
 +\mathcal B_\chi
 \log\frac{\log(2q-1)\log(2q)}{\log(q-1)\log q},\\
\mathcal R_k^\sigma={}&\sum_{a=0}^3\epsilon_a
                         \log\det(I+\Delta_{k,N_a}^\sigma).
\end{aligned}
\]
Since the coefficient of $d_k^\sigma$ is $1+1-1-1=0$,
TWA22 proves
\[
\boxed{F_K(1)+F_B(1)
=\mathcal M_k-\mathcal R_k^\sigma+\delta_k^\sigma,
\qquad
\underline\delta_k^\sigma\leq\delta_k^\sigma
                         \leq\overline\delta_k^\sigma.}
\tag{TWA23}
\]
This is an exact finite correspondence between the FPK coordinates,
the original source comparison and the complete BRU sum.

The original tensor-Gamma reference at BRU $x=0$ has full mass
$(2\pi)^{k/2}$ and density
\[
m_0(u)=\frac{(2\pi)^{k/2}}{c_{k/4}}
                   \frac{|\Gamma(k/4+iu/2)|^2}{2\pi}.
\]
Here the indexed Gamma mass constant is
$c_a=2^{1-2a}\Gamma(2a)$ for $a>0$; the original coordinate
constant remains $c=k/2$.
Repeat TWA18--TWA22 with the original tensor-Gamma monic polynomials
$p_{k,n}(S)=i^n b_n^{(k/4)}((S-c)/i)$ in every kernel entry,
replacing the one-factor polynomials $p_n$ of TWA18, and with
precisely their original FPK.27 norms
\[
\gamma_{k,n}=(2\pi)^{k/2}
                   \frac{n!\Gamma(n+k/2)}{\Gamma(k/2)}
\]
and the full coefficients
\[
C_{\pm,k}(w)=(2\pi)^{-k/4}\sqrt{\Gamma(k/2)}
              \frac{2^{-k/4\pm w/2}}{\Gamma(k/4\pm w/2)}.
\]
The substitution is made everywhere, including in each critical
limiting block. The exact determinant argument is unchanged, and
its exponents remain precisely TWA21. Denote its signed finite
residual by $\mathcal R_k^0$. It follows that
\[
\boxed{\begin{aligned}
F_K(0)+F_B(0)&=\mathcal M_k-\mathcal R_k^0,\\
\Delta_k^\Gamma
&=\delta_k^\sigma+\mathcal R_k^0-\mathcal R_k^\sigma,\\
\mathcal R_k^0-\mathcal R_k^\sigma+\underline\delta_k^\sigma
&\leq\Delta_k^\Gamma
\leq\mathcal R_k^0-\mathcal R_k^\sigma+\overline\delta_k^\sigma.
\end{aligned}}\tag{TWA24}
\]
Thus both Gamma presentations have their exact maps to the original
source, with every mass factor and both signed residuals retained.

\subsection{The exact domain of the fixed-packet limit}

FPK proves $\Delta_{k,N}^\sigma\to0$ and its tensor-Gamma
counterpart as $N\to\infty$ with $k$, every root, every multiplicity
and the entire observation fixed. Its coefficient bounds are
uniform on fixed compact sets and for fixed derivative orders.
For the present packet, $|\Im w_i|$ can reach $k\gamma$, the
number of blocks is $(k+1)^2$, the derivative order can reach
$\ell_k-1$, and the positive limiting matrix has dimension $q$.
In the tensor-Gamma version, its Gamma parameter is also $k/4$.
None of the fixed-parameter convergence statements bounds the
four matrices TWA20 at $N=q-1,q,2q-1,2q$ as $k$ grows.
Both source theorems themselves respect these quantifiers: TW's
window limit follows from its explicit constants, while FPK
explicitly states the absence of packet-uniform error.

The exact residual matrices remain usable at every finite endpoint.
For example let $l_{k,N}$ and $u_{k,N}$ be the smallest and largest
eigenvalues of $I+\Delta_{k,N}^\sigma$. Positivity gives
$0<l_{k,N}\le u_{k,N}$, and multiplying all positive eigenvalues
proves the unconditional finite bound
\[
q\log l_{k,N}\le\log\det(I+\Delta_{k,N}^\sigma)
                         \le q\log u_{k,N}.
\]
The original dimension $q$ and the complete finite matrix are
present. No unproved bound on these eigenvalues is inserted.

There is a valid independent-degree consequence of the fixed-packet
limit. Hold this actual $k,\chi$ fixed and let an integer $n\to\infty$
with $n\ge q$. Apply TWA22 at $n-1,n,2n-1,2n$. The four residual
logarithms tend to zero; their constant terms cancel. The product
of degree ratios tends to four, and the ratio of their logarithm
products tends to one. Therefore
\[
\log\frac{\det G_{n-1}^\sigma\det G_n^\sigma}
               {\det G_{2n-1}^\sigma\det G_{2n}^\sigma}
\longrightarrow2\mathcal L_k\log2.
\tag{TWA25}
\]
The same proof applies to $m_0$ with its full original mass.
TWA25 is the exact fixed-packet conclusion. TWA10--TWA24 give
the original moving-window identities and bounds without
substituting $n=q(k)$ into that independent limit.


\subsection{The actual tensor-Gamma return on the class $k=4l+1$}

The complete finite Christoffel proof GCR1--44 evaluates the
two-reference difference in TWA24 on this tensor class.
We state its full finite matrices and their exact return here.
Let $l=(k-1)/4\geq0$ and retain the original $c=k/2$,
$\ell_k=1+k(m-1)$, $\chi$, all full jets and every unit
coefficient in TWA1. Put
\[
\begin{gathered}
 t_j=2j+\tfrac12,\quad \xi_j=c+t_j,\quad
 f_l(S)=\prod_{j=0}^{l-1}(S-c-t_j),\\
 \beta_k=\frac{(2\pi)^{k/2}}{\sqrt2\,\Gamma(k/2)},\qquad
 dm_{0,k}(y)=\beta_k|f_l(c+iy)|^2\,d\sigma(y).
\end{gathered}\tag{TWA26}
\]
Repeated use of $\Gamma(z+1)=z\Gamma(z)$ gives the factor
$4^{-l}\prod_j(y^2+t_j^2)$, and
$4^lc_{k/4}=\sqrt2\Gamma(k/2)$ gives the stated coefficient.
Since $f_l(c+iy)=i^l\prod_j(y+it_j)$, the equality is in
the original coordinate and retains the respective masses
$(2\pi)^{k/2}$ and $\sqrt{2\pi}$. For $l=0$ the polynomial
and coefficient are both one.

For the same fixed-Gamma monic polynomials $p_n$ and
norms $\gamma_n$ of TWA18, let $\mathcal E_{\chi,\mu}$
be each original divided jet, and define
\[
\begin{gathered}
 (K_N)_{\mu,\nu}
 =\sum_{n=0}^N\mathcal E_{\chi,\mu}(p_n)
      \overline{\mathcal E_{\chi,\nu}(p_n)}/\gamma_n,\\
 (K_{f,L})_{i,j}
 =\sum_{n=0}^L p_n(\xi_i)\overline{p_n(\xi_j)}/\gamma_n,\qquad
 (K_{\chi,f;L})_{\mu,j}
 =\sum_{n=0}^L\mathcal E_{\chi,\mu}(p_n)
                         \overline{p_n(\xi_j)}/\gamma_n,\\
 (U_N)_{\mu,j}=\mathcal E_{\chi,\mu}(p_{N+j})
                               /\sqrt{\gamma_{N+j}}\quad(1\leq j\leq l),\\
 V_N=K_{\chi,f;N+l}K_{f,N+l}^{-1/2},\quad
 W_N=(\,U_N\ \ V_N\,),\quad J_l=\operatorname{diag}(I_l,-I_l),\\
 C_N=K_N+U_NU_N^*-V_NV_N^*>0,\qquad
 d_N^{\rm Chr}=\frac{\det C_N}{\det K_N}
       =\det(I_{2l}+J_lW_N^*K_N^{-1}W_N)>0.
\end{gathered}\tag{TWA27}
\]
These $U_N$ are kernel columns, distinct from the original
full Taylor multiplication unit. All inverses in TWA27
are on positive Grams. Indeed $k$ is odd, so the imaginary
part of every original root is $(2b-k)\gamma\ne0$;
the $\xi_j$ are distinct real roots. Thus their product
jets are onto at degree $N+l\geq q+l-1$. The joint
positive kernel has bottom-right Schur complement
$C_N=K_{\chi,N+l}-K_{\chi,f;N+l}
K_{f,N+l}^{-1}K_{f,\chi;N+l}$, proving positivity.
The finite sum difference
$K_{\chi,N+l}-K_N=U_NU_N^*$ proves its formula.
Taking the determinant of
$\left(\begin{smallmatrix}I&A\\-B&I\end{smallmatrix}\right)$
by either diagonal block proves
$\det(I+AB)=\det(I+BA)$ and hence the last equality.

The entire mixed product in this small determinant is
retained by its positive factorization
\[
\begin{gathered}
 A_N^{\rm Chr}=I_l+U_N^*K_N^{-1}U_N,\qquad
 d_N^{\rm Chr}=\det A_N^{\rm Chr}\det B_N^{\rm Chr},\\
 B_N^{\rm Chr}=I_l-V_N^*K_N^{-1}V_N
 +V_N^*K_N^{-1}U_N(A_N^{\rm Chr})^{-1}U_N^*K_N^{-1}V_N>0.
\end{gathered}\tag{TWA28}
\]
Multiplying the full Woodbury inverse by
$K_N+U_NU_N^*$ proves this identity with its positive
mixed term; positivity is the same Schur complement.

Multiplication by $f_l$ sends the exact remainder fibre
$\{P\in\mathcal P_N:J_NP=z\}$ bijectively onto polynomials
in $\mathcal P_{N+l}$ with product jets
$(0,U_f\mathsf E_\chi z)$, where $U_f$ is full Taylor
multiplication by $f_l$. Its determinant is the unchanged
nonzero product
$\prod_{a,b}\prod_{j=0}^{l-1}
((2a-k)\delta-t_j+i(2b-k)\gamma)^{\ell_k}$.
Vanishing of the $f_l$ jets proves divisibility by $f_l$;
division proves the inverse on this precise fibre.
Taking its minimum with the positive joint kernel gives
$G_N^0=\beta_k\mathsf E_\chi^*U_f^*C_N^{-1}U_f\mathsf E_\chi$.
Thus, before any four-term sum,
\[
 \log\det G_N^0-\log\det G_N^\sigma
      =q\log\beta_k+2\log|\det U_f|-\log d_N^{\rm Chr}.
 \tag{TWA29}
\]
The original observation acts in product jets by
$(0,Y)\mapsto\Lambda\mathsf E_\chi^{-1}U_f^{-1}Y$,
so its full unit and kernel are retained, as are the
relation directions $f_l\chi\mathcal P_{N-q}$.
For a different packet with a common root of orders
$a$ in $f$ and $b$ in $\chi$, the complete GCR6--8
construction uses the product jets of length $a+b$,
first $a$ zero entries and then the invertible
multiplication germ $f/(S-\zeta)^a$ on the next $b$.
It proves the same exact fibre map without inverting
a zero resultant. The present no-collision statement
was proved directly from the actual root coordinates.

Let the four $N_a$ and signs $\epsilon_a$ be exactly those
preceding TWA23. Summing TWA29 cancels its two degree-independent
terms only after retaining their complete values:
\[
\begin{gathered}
 \mathfrak T_k:=F_{0,k}-F_{\sigma,k}
       =-\sum_{a=0}^3\epsilon_a\log d_{N_a}^{\rm Chr},
 \qquad \mathcal R_k^0-\mathcal R_k^\sigma=-\mathfrak T_k,\\
 \Delta_k^\Gamma=\delta_k^\sigma-\mathfrak T_k,\qquad
 -\mathfrak T_k+\underline\delta_k^\sigma
       \leq\Delta_k^\Gamma
       \leq-\mathfrak T_k+\overline\delta_k^\sigma .
\end{gathered}\tag{TWA30}
\]
This evaluates the original tensor-Gamma residual as four
explicit determinants of size $2l$, including all cross entries.
It is the actual value, rather than a discarded reference term.
For $l=0$ the determinants are empty and the return is zero.
The DMR estimate in this body has $k\geq3$, so on this tensor
class its asymptotic statements concern $k\geq5$.

For completeness its finite scalar enclosure also has
fully evaluated moving constants. For $l\geq1$ define
\[
\begin{gathered}
 \mathscr U_{l,N}=\prod_{j=0}^{l-1}[4(N+j+1)^2+t_j^2],
 \quad D_{l,N}^{\rm Chr}
   =\prod_{a=0}^{N}\prod_{r=0}^{2l-1}(a+r+\tfrac12),\\
 X_{l,N}=(N+1)\log\mathscr U_{l,N}-\log D_{l,N}^{\rm Chr}\geq0,
 \qquad Z_k=q\sum_{a=0}^3\epsilon_a\log\mathscr U_{l,N_a},\\
 Z_k-X_{l,q-1}-X_{l,q}\leq\mathfrak T_k
       \leq Z_k+X_{l,2q-1}+X_{l,2q}.
\end{gathered}\tag{TWA31}
\]
The original Gamma Jacobi coefficients are at most their
indices, giving $\|yP\|_\sigma\leq2(d+1)\|P\|_\sigma$
at degree $d$. The identity
$\|(y+it)P\|_\sigma^2=\|yP\|_\sigma^2+t^2\|P\|_\sigma^2$
therefore proves the upper full-vector factor
$\mathscr U_{l,N}$ by successive multiplication.
The ratio of the exact original monic norms is
$\gamma_{k,a}/(\beta_k\gamma_a)
=\prod_{r=0}^{2l-1}(a+r+1/2)$, so multiplying them proves
the determinant ratio $D_{l,N}^{\rm Chr}$ on the same
source dimension $N+1$. The Schur proof of TWA15 then
bounds its quotient deficit between zero and $X_{l,N}$.
Keeping the four signs proves the last line of TWA31.
The complete product estimates GCR34--38, with their
explicit remainders, give
$X_{l,N_a}/[2l(N_a+1)]\to1+\log2$ and
$Z_k/(kq)\to-\log2$. Hence along $k\equiv1\pmod4$,
\[
\begin{gathered}
 -1-2\log2\leq\liminf\frac{\mathfrak T_k}{kq}
       \leq\limsup\frac{\mathfrak T_k}{kq}\leq2+\log2,\\
 -2-\log2\leq\liminf\frac{\Delta_k^\Gamma}{kq}
       \leq\limsup\frac{\Delta_k^\Gamma}{kq}\leq1+2\log2.
\end{gathered}\tag{TWA32}
\]
The second line follows from TWA30 and the proved
$\delta_k^\sigma/(kq)\to0$. These are enclosures, not
limits or a claim of $o(kq)$ for $\mathfrak T_k$.
All other tensor classes retain the literal two-residual
identity TWA24 and the finite interval TWA15.
